% Main claim: update placement trades repeated computation for communication.
\section{Introduction}

Evolution strategies (ES) train a network without backpropagation: perturb the weights
with random noise, score each perturbed model, and move the weights toward the candidates
that scored well. The candidates form a \emph{population}. Their evaluations can run
independently, making ES attractive when many devices are available.

Distributed ES is often described as cheap to communicate because devices exchange only
scalar fitness scores \citep{salimans2017}. Those scores are enough to reconstruct the
update when every device can regenerate the perturbations. But each device must then
repeat the entire reconstruction. An alternative is to reconstruct a partial update on
each device and add the partial updates with an \emph{all-reduce}, an operation that
returns their sum to every device. This sends a model-sized array, but divides the
reconstruction work across devices.

For example, with four devices and 100 candidates, each device evaluates 25 candidates.
After collecting the scores, every device can reconstruct the update from all 100
perturbations, or reconstruct its own 25 contributions and combine the partial updates.
The choice is whether the saved computation is worth the added communication.

The perturbation design changes this balance. \citet{qiu2025} use full-rank Gaussian
noise, regenerated from a seed whenever needed, which saves memory but makes update
reconstruction expensive. EGGROLL \citep{eggroll2025} uses products of thin matrices as
low-rank perturbations. These permit efficient batched evaluation with a shared base
weight matrix and make the update cheap to reconstruct. Neither paper compares update placements, and they do not compare both perturbation
designs in one implementation. Existing systems illustrate two ways of keeping the
complete reconstruction together: Algorithm 2 of \citet{salimans2017} repeats it on
every worker, while EGGROLL's vLLM setup reconstructs on a primary rank and broadcasts
the updated model \citep[Appendix L.2]{eggroll2025}. We test whether dividing that
reconstruction work is worth its communication cost.

We built a JAX \citep{jax2018} library\ifanonymous\else, \textbf{shardes} \citep{shardes},\fi{} to make that
comparison. Every device holds the complete model; the population is what is split
across devices. The library supports both perturbation designs under both placements,
so their costs can be compared with the same model and evaluation code. We ask three
questions:
\begin{itemize}
  \item \textbf{When does splitting the update pay?} A timing model explains the balance
    between repeated computation and communication. On the block benchmark, splitting
    wins for full-rank noise on a single host, while replication wins for rank-1 noise
    at the larger matrix width.
  \item \textbf{Does the result carry over?} The placement preference holds on
    Qwen2.5-0.5B on the TPU. Across a slow connection between hosts, replication wins
    for large full-rank updates too.
  \item \textbf{Do cheaper perturbations remain useful?} On matched block configurations,
    rank-1 generation time is 0.22 times the dense time on the A100 and 0.08 times on
    the v5e. The tested ranks have close one-update alignment and reach similar observed
    final Countdown accuracy, but differ early in training.
\end{itemize}
The placement measurements are the main result. The learning experiments check the
practical relevance of the cheaper perturbations within one model family and one task.
